\documentclass[12pt]{article}
\usepackage{openwork}
\usepackage{amssymb}
\usepackage[english,russian]{babel}
\usepackage{cite}
\usepackage[hidelinks]{hyperref} 

\title{Analysis of the Relationship Between Electroencephalography (EEG) and
Functional Magnetic Resonance Imaging (fMRI) Signals}
\author[1]{Р.И Шевчук}
\author[2]{А.В Грабовой}
\author[3]{Д.Д Дорин}
\affil[1]{shevchuk.r@phystech.edu}
\affil[2]{grabovoy.av@phystech.edu}
\affil[3]{dorin.dd@phystech.edu}
\date{}

\begin{document}
\hypersetup{final}
\maketitle

\begin{abstract}
Одновременные записи ЭЭГ-фМРТ позволяют фиксировать реакции мозга (как слабые, так и сильные) на различные стимулы, например просмотр видео или движения пальцев. ЭЭГ регистрирует электрическую активность мозга на уровне кожи головы, тогда как фМРТ отслеживает гемодинамические изменения, прежде всего изменения насыщенности крови кислородом в областях мозга. Ряд исследований указывает на то, что эти два метода являются взаимодополняющими.

Цель данной статьи - проанализировать взаимосвязь между ЭЭГ и фМРТ путём кодирования данных обоих методов в общее латентное пространство, а также применить статистические тесты для исследования их взаимной зависимости.

\textbf{Ключевые слова:} Мозг, ЭЭМ, фМРТ, contrastive learning, neural network
\end{abstract}

\section{Введение}
Электроэнцефалография (ЭЭГ) и функциональная магнитно-резонансная томография (фМРТ) являются одними из наиболее широко используемых неинвазивных методов исследования человеческого мозга, поскольку предоставляют комплементарные представления нейронной активности. ЭЭГ напрямую отражает электрические процессы в нейронных популяциях с миллисекундным временным разрешением, тогда как фМРТ измеряет медленный гемодинамический отклик, связанный с метаболической активностью, обеспечивая высокое пространственное разрешение.

Одним из ключевых открытых вопросов остаётся природа взаимосвязи между ЭЭГ и фМРТ. С одной стороны, различия во временных масштабах (от миллисекунд у ЭЭГ до секунд у фМРТ), а также различие биофизических механизмов регистрации сигналов позволяют предположить, что эти модальности отражают принципиально разные аспекты нейронной динамики. С другой стороны, накопленные эмпирические данные указывают на существование между ними систематической связи.

В частности, было показано, что связь между функциональной связностью, измеренной с помощью фМРТ, и функциональной связностью, полученной из М/ЭЭГ, во многом определяется структурной связностью белого вещества мозга, которую восстанавливают на основе данных диффузионной МРТ(dMRI) \cite{Sadaghiani2020}. Кроме того, было продемонстрировано, что кросс-модальная связность может быть разложена на две компоненты - общую и дополнительную, которая различается между двумя модальностями\cite{Wirsich2021}. Это, в совокупности, указывает на наличие нетривиальной и многоуровневой взаимосвязи между ЭЭГ и фМРТ.
%~\cite{wirsich_eeg_fmri}.

Особый интерес представляет анализ одновременных записей ЭЭГ-фМРТ, поскольку они позволяют напрямую исследовать связь между электрофизиологическими процессами и гемодинамическим откликом в одних и тех же временных интервалах.
Ранние исследования одновременной ЭЭГ-фМРТ в основном были сосредоточены на так называемом EEG-informed fMRI подходе, при котором временные характеристики ЭЭГ (например, мощность в частотных диапазонах или событие-связанные потенциалы) использовались в качестве регрессоров для анализа BOLD-сигнала. Эти работы продемонстрировали, что флуктуации спектральной мощности ЭЭГ и отдельные электрофизиологические события статистически связаны с региональной активацией фМРТ, однако такой подход в первую очередь описывает локальные или сетевые активации и не даёт ответа на вопрос о согласованности функциональной связности между модальностями.

Более поздние исследования сместили фокус к анализу динамики, связывая временные изменения ЭЭГ-показателей с динамической функциональной связностью фМРТ. В рамках этого направления было показано, что изменения спектральной мощности ЭЭГ и характеристики микросостояний могут сопутствовать переходам между различными состояниями фМРТ-коннектома. Эти результаты указывают на возможную связь между электрофизиологической активностью на быстрых временных масштабах и медленными перестройками функциональных сетей, однако остаются чувствительными к выбору временных окон, метрик связности и пространственного представления ЭЭГ.

Отдельную линию исследований составляет прямое сравнение whole-brain функциональной связности, полученной из ЭЭГ и фМРТ. В таких работах строятся коннектомы для обеих модальностей с использованием единой или сопоставимой парцелляции мозга, после чего проводится сопоставление матриц связности. Было показано, что коннектомы ЭЭГ и фМРТ демонстрируют умеренную, но статистически значимую корреляцию, воспроизводимую между различными частотными диапазонами ЭЭГ и различными экспериментальными установками. Особенно важным является то, что эта связь сохраняется при использовании данных одновременной записи и воспроизводится на разных магнитных полях и конфигурациях ЭЭГ, что указывает на существование общей функциональной основы нейронной организации.

Тем не менее существующие подходы в основном опираются на корреляционные меры и заранее заданные признаки (мощность, фазовая или амплитудная связность), а также на усреднение по времени или по группам испытуемых. Это существенно ограничивает интерпретацию результатов и затрудняет переход от описания статистического сходства к пониманию латентных нейронных источников, лежащих в основе наблюдаемых сигналов. Кроме того, проблема обратной задачи ЭЭГ и влияние артефактов одновременной записи остаются ключевыми факторами, искажающими оценку межмодальной связности.

В данной работе предлагается рассматривать взаимосвязь ЭЭГ и фМРТ как задачу восстановления общего латентного пространства нейронной активности, в котором обе модальности являются различными проекциями одного и того же скрытой динамической системы. Основная гипотеза состоит в том, что пространственно и временно локализованные изменения BOLD-сигнала статистически связаны с нейронными источниками ЭЭГ, и что эта связь может быть выявлена с помощью многомерных и нелинейных методов анализа временных рядов. Для этого используются методы канонического корреляционного анализа и реконструкции латентного пространства, позволяющие выйти за рамки парных корреляций. 

\section{Постановка задачи}

Рассмотрим одновременную запись двух модальностей: электроэнцефалографии (EEG) и функциональной МРТ (fMRI), измеренных в течение одного и того же эксперимента/сессии.
Пусть
\[
\mathbf{X}^{(e)} \in \mathbb{R}^{C' \times T'}, 
\qquad
\mathbf{X}^{(f)} \in \mathbb{R}^{C \times T},
\]
где $\mathbf{X}^{(e)}$ - EEG-сигналы, $\mathbf{X}^{(f)}$ - fMRI-сигналы.
Здесь $C'$ и $C$ обозначают число пространственных каналов (например сенсоров EEG, а также число regions of interest в выбранной парцелляции fMRI), а $T'$ и $T$ - число временных отсчётов.
Как правило, $T' > T$, поскольку EEG имеет существенно более высокую временную дискретизацию по сравнению с fMRI.

В общем случае $C' \neq C$ из-за различий в аппаратных конфигурациях и/или различий в пространственной дискретизации (например, разные атласы и парцелляции для fMRI, либо разные способы приведения EEG к ROI/источникам).
Следовательно, требуется метод, который:
\begin{enumerate}
    \item устойчив к несовпадению пространственных размерностей $C' \neq C$
    \item учитывает различие временных масштабов $T' \neq T$
\end{enumerate}

\subsection{Модель представления}

Обучаются два энкодера, каждый для своей модальности:
\[
E_e:\mathbb{R}^{C' \times T'} \to \mathbb{R}^{1 \times T},
\qquad
E_f:\mathbb{R}^{C \times T} \to \mathbb{R}^{1 \times T},
\]
которые отображают входные многоканальные временные ряды в \emph{одномерные} временные ряды
\[
\mathbf{z}^{(e)} = E_e(\mathbf{X}^{(e)}) \in \mathbb{R}^{T},
\qquad
\mathbf{z}^{(f)} = E_f(\mathbf{X}^{(f)}) \in \mathbb{R}^{T}.
\]

Далее вводим нормировку для стабилизации контрастивного обучения:
\[
\tilde{\mathbf{z}}^{(e)} = \frac{\mathbf{z}^{(e)}}{\|\mathbf{z}^{(e)}\|_2}, 
\qquad
\tilde{\mathbf{z}}^{(f)} = \frac{\mathbf{z}^{(f)}}{\|\mathbf{z}^{(f)}\|_2}.
\]

\subsection{Обучение с clip-like loss}

Пусть в батче присутствуют $N$ пар одновременных наблюдений
$\{(\mathbf{X}^{(e)}_i,\mathbf{X}^{(f)}_i)\}_{i=1}^N$,
где пара с одинаковым индексом $i$ является \emph{положительной} (positive), а пары с различными индексами $(i,j)$, $i\neq j$, являются \emph{отрицательными} (negative).

Определим матрицу сходства между латентными представлениями как скалярные произведения:
\[
s_{ij} = \left\langle \tilde{\mathbf{z}}^{(e)}_i,\, \tilde{\mathbf{z}}^{(f)}_j \right\rangle
\in [-1,1].
\]
Введём температурный параметр $\tau > 0$.

Тогда clip-like контрастивный лосс задаётся:
\[
\mathcal{L} = \frac{1}{2}\left(\mathcal{L}_{e\to f} + \mathcal{L}_{f\to e}\right),
\]
где
\[
\mathcal{L}_{e\to f} 
= -\frac{1}{N}\sum_{i=1}^{N}\log
\frac{\exp\left(s_{ii}/\tau\right)}
{\sum_{j=1}^{N}\exp\left(s_{ij}/\tau\right)},
\]
\[
\mathcal{L}_{f\to e} 
= -\frac{1}{N}\sum_{i=1}^{N}\log
\frac{\exp\left(s_{ii}/\tau\right)}
{\sum_{j=1}^{N}\exp\left(s_{ji}/\tau\right)}.
\]
Минимизация $\mathcal{L}$ приводит к тому, что латентные одномерные временные ряды $\mathbf{z}^{(e)}_i$ и $\mathbf{z}^{(f)}_i$, соответствующие одной и той же записи, становятся максимально согласованными, в то время как несоответствующие пары отталкиваются.

\bibliographystyle{plain}
\begin{thebibliography}{9}
\bibitem{Sadaghiani2020} Sadaghiani, S., & Wirsich, J. (2020). Intrinsic connectome
organization across temporal scales:
New insights from cross-modal
approaches. Network Neuroscience,
4(1), 1–29. \hyperlink{https://doi.org/10.1162/netn_a_00114}{Article}}

\bibitem{Wirsich2021}
Jonathan Wirsich, João Jorge, Giannina Rita Iannotti, Elhum A Shamshiri, Frédéric Grouiller, Rodolfo Abreu, François Lazeyras, Anne-Lise Giraud, Rolf Gruetter, Sepideh Sadaghiani, Serge Vulliémoz (2021) The relationship between EEG and fMRI connectomes is reproducible across simultaneous EEG-fMRI studies from 1.5 T to 7T \hyperlink{https://www.sciencedirect.com/science/article/pii/S1053811921001415}{Article}

\end{thebibliography}

\end{document}